\chapter{Single-Source Shortest Paths}

You are standing at a street corner in a city you have never visited, holding a
phone that promises to route you to a friend's apartment across town. In the
half-second before the blue line appears, something remarkable happens: a piece
of software considers an astronomical number of possible walking routes -- left
here, right there, double back, cut through the park -- and returns the single
shortest one. Not a pretty-good route. The shortest. And it does this not by
trying every route (there are far too many) but by exploiting fundamental
structural properties that govern how short paths are built out of shorter ones.

This chapter is about those facts and the algorithms that ride on them. We are
given a graph whose edges carry \emph{weights} -- think distances, times, costs,
tolls -- and a single \emph{source} vertex. We want the cheapest path from the
source to every other vertex. That is the \emph{single-source shortest-paths}
problem, and it is one of the most useful problems in all of computing: routing,
networking, scheduling, currency arbitrage, and a hundred other things reduce to
it.

We take it slowly. We start by pinning down what a shortest path even
\emph{is}, and uncover the structural facts that make the whole enterprise
possible. Then comes the one operation -- \emph{relaxation} -- that every
algorithm in this chapter is secretly built from. With that in hand we study
three algorithms in increasing order of cleverness: Bellman--Ford (slow but
fearless, handles negative edges and detects the pathological cases), the DAG
algorithm (blisteringly fast when the graph has no cycles), and Dijkstra's
algorithm (the elegant workhorse behind your phone's blue line). Along the way we
will prove, carefully but without ceremony, that these algorithms actually return
shortest paths.

\chapterline

\section{What is a shortest path, exactly?}

Let us set up the notation carefully, because sloppiness here causes endless
confusion later.

We have a directed graph $G = (V, E)$ with $|V| = n$ vertices and $|E| = m$
edges. (We write $n$ and $m$ throughout; CLRS often writes $|V|$ and $|E|$.) Each
edge $(u, v) \in E$ has a real-valued \emph{weight} $w(u, v)$. A \emph{path} from
$u$ to $v$ is a sequence of vertices $\langle v_0, v_1, \dots, v_k\rangle$ with
$v_0 = u$, $v_k = v$, and every consecutive pair $(v_{i-1}, v_i)$ an edge of $G$.
The \emph{weight} of that path is just the sum of its edge weights:
\[
w(p) \;=\; \sum_{i=1}^{k} w(v_{i-1}, v_i).
\]

\begin{definition}[Shortest-path weight]\label{ch24:def-delta}
The \emph{shortest-path weight} from $u$ to $v$, written $\delta(u, v)$, is
\[
\delta(u, v) =
\begin{cases}
\min\{\, w(p) : p \text{ is a path from } u \text{ to } v \,\} & \text{if such a path exists,}\\
\infty & \text{otherwise.}
\end{cases}
\]
A \emph{shortest path} from $u$ to $v$ is any path $p$ with $w(p) = \delta(u, v)$.
\end{definition}

Notice the word \emph{any}. There can be many shortest paths tied for the
minimum, and that is fine; we want one of them (and its weight). And since
$\delta$ is a \emph{minimum}, a minimum over an infinite set of paths only
makes sense if that set is bounded below. Hold that thought -- it is exactly
where negative-weight cycles will bite us.

\begin{intuitionbox}
Picture the graph as a road network. The vertices are intersections, the
directed edges are one-way street segments, and the weight on an edge is how many
minutes that segment takes. A path is a legal driving route; its weight is the
total driving time. $\delta(u,v)$ is the fastest possible time to get from $u$ to
$v$. Our whole job, given a starting intersection $s$, is to fill in a little
table: ``fastest time from $s$ to here'' for every intersection.
\end{intuitionbox}

\subsection{Why negative edges, and why negative cycles are poison}

Weights are not always distances, and they are not always positive. If an edge
weight is a financial gain or loss, or an elevation change, it can be negative.
Negative \emph{edges} are perfectly fine and meaningful. The trouble is negative
\emph{cycles}.

Suppose there is a cycle -- a path that returns to where it started -- whose total
weight is negative. Then for any vertex on that cycle, you can lower your path
weight by looping around once more, and again, and again, forever. The minimum in
Definition~\ref{ch24:def-delta} is then $-\infty$: there is no shortest path. We
say $\delta(u,v) = -\infty$ for any $v$ reachable from such a cycle (reachable
from $u$).

\begin{trapbox}
A negative \emph{edge} is not a problem. A negative-weight \emph{cycle that is
reachable from the source} is. The two get confused constantly. Dijkstra's
algorithm forbids negative edges entirely (we will see precisely why). Bellman--Ford
tolerates negative edges and, as a bonus, \emph{detects} negative cycles and tells
you they exist instead of silently returning garbage.
\end{trapbox}

There is also the gentler cousin: a \emph{positive} or \emph{zero}-weight cycle.
Those never help -- looping around a cycle of nonnegative weight can only keep
your weight the same or make it worse. So we can always assume a shortest path
has \emph{no repeated vertices}: it is a \emph{simple} path, with at most $n$
vertices and hence at most $n - 1$ edges. Remember that number, $n-1$. It will
reappear as the magic iteration count in Bellman--Ford.

\section{The two structural facts everything rests on}

Before any algorithm, we extract the facts that make shortest paths tractable.
There are really only two, and both are intuitive once you see them.

\subsection{Optimal substructure: shortest paths contain shortest paths}

\begin{intuitionbox}
Consider the natural substructure of optimal routing: suppose the fastest route
from your home to the airport happens to pass through a particular gas station. Then the
piece of that route from home to the gas station must itself be the fastest route
from home to the gas station. Why? Because if there were a faster way to reach the
gas station, you could splice it in and beat your supposedly-fastest airport
route -- contradiction. \emph{Subpaths of shortest paths are shortest paths.}
\end{intuitionbox}

\begin{lemma}[Optimal substructure]\label{ch24:lemma-substructure}
Let $p = \langle v_0, v_1, \dots, v_k\rangle$ be a shortest path from $v_0$ to
$v_k$. For any $0 \le i \le j \le k$, the subpath
$p_{ij} = \langle v_i, \dots, v_j\rangle$ is a shortest path from $v_i$ to $v_j$.
\end{lemma}

\begin{formalbox}
\emph{Proof (cut-and-paste).} Decompose $p$ into three pieces:
$v_0 \xrightarrow{p_{0i}} v_i \xrightarrow{p_{ij}} v_j \xrightarrow{p_{jk}} v_k$,
so $w(p) = w(p_{0i}) + w(p_{ij}) + w(p_{jk})$. Suppose, for contradiction, that
$p_{ij}$ is \emph{not} a shortest path from $v_i$ to $v_j$; then there is a path
$p'_{ij}$ from $v_i$ to $v_j$ with $w(p'_{ij}) < w(p_{ij})$. Replace $p_{ij}$ by
$p'_{ij}$ inside $p$ to get a new $v_0$-to-$v_k$ path of weight
$w(p_{0i}) + w(p'_{ij}) + w(p_{jk}) < w(p)$. That contradicts $p$ being shortest.
Hence $p_{ij}$ is shortest. \qed
\end{formalbox}

This optimal substructure is exactly what licenses dynamic programming and greedy
strategies on shortest paths. It is the same shape of argument you saw in
Chapter~15.

\subsection{The triangle inequality}

The second fact is the shortest-path version of ``a detour is never shorter than
going straight.''

\begin{lemma}[Triangle inequality]\label{ch24:lemma-triangle}
For any edge $(u, v) \in E$ and any source $s$,
\[
\delta(s, v) \;\le\; \delta(s, u) + w(u, v).
\]
\end{lemma}

\begin{formalbox}
\emph{Proof.} One particular way to get from $s$ to $v$ is: take a shortest path
from $s$ to $u$ (weight $\delta(s,u)$), then traverse the single edge $(u,v)$
(weight $w(u,v)$). This is \emph{a} path to $v$, of weight $\delta(s,u)+w(u,v)$.
The shortest-path weight $\delta(s,v)$ is the minimum over \emph{all} paths to
$v$, so it cannot exceed this particular one. (If $\delta(s,u) = \infty$ the
inequality holds directly; if a negative cycle makes $\delta(s,u) = -\infty$, then
$\delta(s,v) = -\infty$ too and again it holds.) \qed
\end{formalbox}

\begin{focusbox}
Two facts. (1) \textbf{Optimal substructure}: chop a shortest path anywhere and
each piece is itself shortest. (2) \textbf{Triangle inequality}:
$\delta(s,v) \le \delta(s,u) + w(u,v)$ for every edge. The first tells us shortest
paths can be \emph{built up} from smaller shortest paths. The second is the
yardstick our algorithms will constantly check estimates against. Everything else
in this chapter is bookkeeping on top of these two.
\end{focusbox}

\section{Relaxation: the one move to rule them all}

Here is the secret. \emph{Every} algorithm in this chapter -- Bellman--Ford, the
DAG algorithm, Dijkstra -- is the same single operation applied to edges, over and
over, in different orders. That operation is called \emph{relaxation}. Understand
it, and the three algorithms become three ways of scheduling it.

We keep two arrays. For each vertex $v$:
\begin{itemize}
\item $d[v]$ -- our current \emph{estimate} of $\delta(s, v)$. It starts at
$\infty$ (``no idea, assume unreachable'') and only ever goes \emph{down}. It is
always an \emph{upper bound}: $d[v] \ge \delta(s,v)$ at all times.
\item $\pi[v]$ -- the \emph{predecessor} of $v$ on the best path we have found so
far. Following $\pi$ pointers backward from $v$ reconstructs the path. Starts at
\textsc{nil}.
\end{itemize}

Initialization sets everyone to ``unknown'' except the source, which is zero
distance from itself.

\begin{algorithm}[H]
\caption{\textsc{Initialize-Single-Source}($G, s$)}
\KwIn{Graph $G=(V,E)$, source vertex $s$.}
\KwOut{Arrays $d$ and $\pi$ initialized.}
\For{each vertex $v \in V$}{
  $d[v] \leftarrow \infty$\;
  $\pi[v] \leftarrow \textsc{nil}$\;
}
$d[s] \leftarrow 0$\;
\end{algorithm}

Now the star of the show. To \emph{relax} an edge $(u,v)$ is to ask one question:
``Can I reach $v$ more cheaply by going through $u$?'' If yes, we improve the
estimate.

\begin{algorithm}[H]
\caption{\textsc{Relax}($u, v, w$)}
\KwIn{Edge $(u,v)$ with weight $w(u,v)$; current $d$, $\pi$.}
\KwOut{$d[v]$ and $\pi[v]$ possibly improved.}
\If{$d[v] > d[u] + w(u,v)$}{
  $d[v] \leftarrow d[u] + w(u,v)$ \tcp*{found a cheaper route to $v$ via $u$}
  $\pi[v] \leftarrow u$\;
}
\end{algorithm}

\begin{intuitionbox}
Relaxation is a rumor spreading through the graph. Vertex $u$ currently believes
it can be reached from $s$ in $d[u]$ time. It turns to its neighbor $v$ and says,
``I can get to you in $d[u] + w(u,v)$. Is that better than your current best?'' If
so, $v$ updates its belief and remembers that the good news came from $u$. The
word ``relax'' comes from physics/optimization: we are loosening an
over-tight upper bound, letting it slip down toward the truth.
\end{intuitionbox}

Why does this single move suffice? Because of the triangle inequality. The true
distances satisfy $\delta(s,v) \le \delta(s,u) + w(u,v)$ for every edge. An
algorithm has finished only when its estimates likewise satisfy
$d[v] \le d[u] + w(u,v)$ for \emph{every} edge -- when no edge can be relaxed any
further. At that point, as we will prove, the estimates have hit the truth. So
the entire game is: relax edges, in some clever order, until none can be improved.

Two small lemmas make this precise. They are the workhorses behind every
correctness proof to come, so let us state them plainly.

\begin{lemma}[Relaxation never undershoots the truth]\label{ch24:lemma-upperbound}
If we run \textsc{Initialize-Single-Source} and then perform any sequence of
relaxation steps, then $d[v] \ge \delta(s,v)$ holds for every vertex $v$ at all
times. Moreover, once $d[v]$ reaches $\delta(s,v)$, it never changes again.
\end{lemma}

\begin{formalbox}
\emph{Proof sketch (induction on the number of relaxations).} Initially
$d[s]=0=\delta(s,s)$ and $d[v]=\infty\ge\delta(s,v)$ for $v\ne s$. Now suppose the
invariant holds and we relax $(u,v)$, possibly setting
$d[v] = d[u] + w(u,v)$. By the inductive hypothesis $d[u]\ge\delta(s,u)$, so
$d[u]+w(u,v) \ge \delta(s,u)+w(u,v) \ge \delta(s,v)$, the last step by the
triangle inequality (Lemma~\ref{ch24:lemma-triangle}). Thus the new $d[v]$ is
still $\ge\delta(s,v)$. Since $d[v]$ can never drop below $\delta(s,v)$ and relax
only ever decreases it, once it equals $\delta(s,v)$ it is stuck there. \qed
\end{formalbox}

\begin{lemma}[Convergence]\label{ch24:lemma-convergence}
Let $s \rightsquigarrow u \to v$ be a shortest path to $v$, so that
$\delta(s,v) = \delta(s,u) + w(u,v)$. Suppose at some moment $d[u] = \delta(s,u)$,
and at some later moment we relax the edge $(u,v)$. Then from that point on
$d[v] = \delta(s,v)$.
\end{lemma}

\begin{formalbox}
\emph{Proof.} By Lemma~\ref{ch24:lemma-upperbound}, once $d[u]=\delta(s,u)$ it
stays there, so at the moment we relax $(u,v)$ we have $d[u]=\delta(s,u)$.
Relaxation guarantees afterward $d[v] \le d[u]+w(u,v) = \delta(s,u)+w(u,v) =
\delta(s,v)$. But also $d[v]\ge\delta(s,v)$ always. Hence $d[v]=\delta(s,v)$, and
it cannot change afterward. \qed
\end{formalbox}

The convergence lemma is the lever every algorithm pulls. It says: \emph{if you
relax the edges of a shortest path in order, the estimate at the end of that path
becomes correct.} Each algorithm is just a strategy for making sure the edges of
every shortest path get relaxed in the right order.

\begin{walkbox}
Before reading on, take a walk with this question. The convergence lemma needs the
edges of a shortest path relaxed \emph{in path order} -- $d[u]$ correct first,
\emph{then} relax $(u,v)$. If you relax edges in a bad order, you might relax
$(u,v)$ while $d[u]$ is still too big, miss the improvement, and never revisit it.
How could you guarantee good order? Two ideas: relax \emph{everything} many times
so order stops mattering (Bellman--Ford), or process vertices in an order where
$d$-values are correct by the time you use them (topological order for DAGs,
increasing-distance order for Dijkstra). The whole chapter is those two ideas.
\end{walkbox}

\section{Bellman--Ford: relax everything, repeatedly}

The most fearless idea is the simplest. If we are not clever enough to figure out
a good relaxation order, we can be \emph{thorough} instead: relax \emph{every}
edge, then do it again, and again, enough times that no order could have mattered.
How many times is enough? Exactly $n-1$, and we can see why from the convergence
lemma.

A shortest (simple) path has at most $n-1$ edges. Walk along such a path
$s = v_0 \to v_1 \to \dots \to v_k = v$ with $k \le n-1$. After the first full
pass over all edges, the edge $(v_0,v_1)$ certainly got relaxed at some point, and
since $d[v_0]=\delta(s,v_0)=0$ from the start, the convergence lemma makes
$d[v_1]=\delta(s,v_1)$. After the second pass, $(v_1,v_2)$ got relaxed with $d[v_1]$
already correct, so $d[v_2]$ becomes correct. In general, after pass $i$, the
first $i$ vertices on the path have correct $d$-values. So after $k \le n-1$
passes, $d[v]=\delta(s,v)$. Since this holds for every vertex's shortest path
simultaneously, $n-1$ passes settle everyone.

\begin{algorithm}[H]
\caption{\textsc{Bellman-Ford}($G, w, s$)}
\KwIn{Weighted digraph $G=(V,E)$ with weights $w$ (may be negative), source $s$.}
\KwOut{\textsc{true} with correct $d,\pi$ if no negative cycle is reachable; else \textsc{false}.}
\textsc{Initialize-Single-Source}($G,s$)\;
\For{$i \leftarrow 1$ \KwTo $n-1$}{
  \For{each edge $(u,v) \in E$}{
    \textsc{Relax}($u, v, w$)\;
  }
}
\For{each edge $(u,v) \in E$}{
  \If{$d[v] > d[u] + w(u,v)$}{
    \Return \textsc{false} \tcp*{an edge still improves: negative cycle!}
  }
}
\Return \textsc{true}\;
\end{algorithm}

\subsection{The negative-cycle detector, explained}

Look at that final loop. After $n-1$ full passes, if the graph had no reachable
negative cycle, \emph{every} edge would satisfy $d[v] \le d[u]+w(u,v)$ -- nothing
left to relax. So if even one edge still says ``I can improve you,''
$d[v] > d[u]+w(u,v)$, the only explanation is a negative cycle: the relaxations
never converged because there is no finite minimum. The check is a single extra
pass, and it turns a fragile algorithm into a robust one that \emph{tells you when
the question has no answer}.

\begin{formalbox}
\emph{Why the detector is correct.} Suppose a negative-weight cycle
$c = \langle v_0, v_1, \dots, v_k = v_0\rangle$ reachable from $s$ existed and the
algorithm wrongly returned \textsc{true}. Returning \textsc{true} means every edge
on $c$ passed the test: $d[v_i] \le d[v_{i-1}] + w(v_{i-1},v_i)$ for
$i=1,\dots,k$. Sum all $k$ of these inequalities:
\[
\sum_{i=1}^{k} d[v_i] \;\le\; \sum_{i=1}^{k} d[v_{i-1}] \;+\; \sum_{i=1}^{k} w(v_{i-1},v_i).
\]
Because $c$ is a cycle, the vertices on the left and the vertices in the first sum
on the right are the same multiset, so those two sums are equal and cancel. We are
left with $0 \le \sum_{i=1}^{k} w(v_{i-1},v_i) = w(c)$. But $c$ is a negative
cycle, $w(c) < 0$, a contradiction. So a reachable negative cycle always trips the
detector. \qed
\end{formalbox}

\subsection{A complete numeric trace}

Let us run Bellman--Ford on a small graph with a negative edge but \emph{no}
negative cycle, so we can watch relaxation grind its way to the answer. The
vertices are
$s, t, x, y, z$ and the edges (with weights) are:
\[
\begin{aligned}
&(s,t)=6,\ (s,y)=7,\ (t,x)=5,\ (t,y)=8,\ (t,z)=-4,\\
&(x,t)=-2,\ (y,x)=-3,\ (y,z)=9,\ (z,s)=2,\ (z,x)=7.
\end{aligned}
\]

\begin{figure}[h]
\centering
\begin{tikzpicture}[->,>=Stealth,node distance=2.4cm,
  every node/.style={circle,draw,minimum size=7mm,inner sep=1pt,font=\small},
  every edge/.style={draw,font=\scriptsize}]
  \node[fill=sourceblue!25] (s) at (0,1.4) {$s$};
  \node (t) at (2.6,2.8) {$t$};
  \node (x) at (5.2,2.8) {$x$};
  \node (y) at (2.6,0) {$y$};
  \node (z) at (5.2,0) {$z$};
  \draw (s) edge node[above]{$6$} (t);
  \draw (s) edge node[above left]{$7$} (y);
  \draw (t) edge node[above]{$5$} (x);
  \draw (t) edge[bend left=10] node[left]{$8$} (y);
  \draw (t) edge node[pos=0.4,above right]{$-4$} (z);
  \draw (x) edge[bend left=10] node[right]{$-2$} (t);
  \draw (y) edge node[below]{$-3$} (x);
  \draw (y) edge node[below]{$9$} (z);
  \draw (z) edge[bend left=25] node[below left]{$2$} (s);
  \draw (z) edge[bend left=10] node[below]{$7$} (x);
\end{tikzpicture}
\caption{The Bellman--Ford example graph. Source $s$ is shaded. Note the negative
weights on $(t,z)$, $(x,t)$, $(y,x)$, but no negative cycle.}
\label{ch24:bf-graph}
\end{figure}

We fix the edge-relaxation order once and reuse it every pass:
\[
(t,x),\ (t,y),\ (t,z),\ (x,t),\ (y,x),\ (y,z),\ (z,x),\ (z,s),\ (s,t),\ (s,y).
\]
Start with $d[s]=0$ and everyone else $\infty$.

\begin{labbox}
Trace each pass with me. In a pass, walk the edge list above and apply
\textsc{Relax}. An edge does nothing if its tail $u$ still has $d[u]=\infty$.

\medskip
\textbf{Pass 1.} Most edges have an infinite tail and do nothing. The two edges
out of $s$ fire: $(s,t)$ sets $d[t]=0+6=6$; $(s,y)$ sets $d[y]=0+7=7$. Everyone
else is still $\infty$.
\[
d = (s{:}0,\ t{:}6,\ x{:}\infty,\ y{:}7,\ z{:}\infty).
\]

\medskip
\textbf{Pass 2.} Now $t$ and $y$ are finite. $(t,x)$: $d[x]=6+5=11$. $(t,z)$:
$d[z]=6+(-4)=2$. $(y,x)$: $7+(-3)=4 < 11$, so $d[x]=4$. $(y,z)$:
$7+9=16$, not better than $2$. After the pass:
\[
d = (s{:}0,\ t{:}6,\ x{:}4,\ y{:}7,\ z{:}2).
\]

\medskip
\textbf{Pass 3.} $(x,t)$: $4+(-2)=2 < 6$, so $d[t]=2$. Other edges find no
improvement this pass.
\[
d = (s{:}0,\ t{:}2,\ x{:}4,\ y{:}7,\ z{:}2).
\]

\medskip
\textbf{Pass 4.} With $d[t]=2$ now, $(t,z)$: $2+(-4)=-2 < 2$, so $d[z]=-2$. Nothing
else improves.
\[
d = (s{:}0,\ t{:}2,\ x{:}4,\ y{:}7,\ z{:}-2).
\]
That is $n-1 = 4$ passes for $n=5$ vertices. The final pass over all edges finds
\emph{no} improvable edge, so the detector returns \textsc{true}. These are the
true shortest-path weights: $\delta(s,t)=2$, $\delta(s,x)=4$, $\delta(s,y)=7$,
$\delta(s,z)=-2$.
\end{labbox}

Following $\pi$ pointers, the shortest path to $z$ is
$s \to y \to x \to t \to z$ with weight $7-3-2-4=-2$. Pause on that: the
\emph{shortest} route to $z$ goes the ``long way'' through four edges, exploiting
the negative weights. A greedy nearest-first strategy would have grabbed
$s\to t\to z$ for weight $2$ and missed it. This is exactly why Dijkstra cannot
handle negative edges, a point we return to once we meet it.

In this particular example the final predecessor pointers happen to line up into a
single chain: $\pi[y]=s$, $\pi[x]=y$, $\pi[t]=x$, $\pi[z]=t$. That chain
\emph{is} the shortest-paths tree (Figure~\ref{ch24:bf-tree}); reading it backward
from any vertex recovers that vertex's shortest path from $s$.

\begin{figure}[h]
\centering
\begin{tikzpicture}[->,>=Stealth,
  every node/.style={circle,draw,minimum size=7mm,inner sep=1pt,font=\small}]
  \node[fill=sourceblue!25] (s) at (0,0) {$s$};
  \node (y) at (2,0) {$y$};
  \node (x) at (4,0) {$x$};
  \node (t) at (6,0) {$t$};
  \node (z) at (8,0) {$z$};
  \draw[targetgreen,very thick] (s) -- node[above,font=\scriptsize,draw=none,fill=none]{$7$} (y);
  \draw[targetgreen,very thick] (y) -- node[above,font=\scriptsize,draw=none,fill=none]{$-3$} (x);
  \draw[targetgreen,very thick] (x) -- node[above,font=\scriptsize,draw=none,fill=none]{$-2$} (t);
  \draw[targetgreen,very thick] (t) -- node[above,font=\scriptsize,draw=none,fill=none]{$-4$} (z);
  \node[draw=none,fill=none,font=\scriptsize] at (0,-0.7) {$d{=}0$};
  \node[draw=none,fill=none,font=\scriptsize] at (2,-0.7) {$d{=}7$};
  \node[draw=none,fill=none,font=\scriptsize] at (4,-0.7) {$d{=}4$};
  \node[draw=none,fill=none,font=\scriptsize] at (6,-0.7) {$d{=}2$};
  \node[draw=none,fill=none,font=\scriptsize] at (8,-0.7) {$d{=}{-}2$};
\end{tikzpicture}
\caption{The shortest-paths tree the $\pi$ pointers encode after Bellman--Ford
finishes on Figure~\ref{ch24:bf-graph}. Here it degenerates to a single path;
following the arrows from $s$ gives each vertex its $\delta$-value (shown below).}
\label{ch24:bf-tree}
\end{figure}

\begin{complexitybox}
\textbf{Bellman--Ford.} The main loop runs $n-1$ passes, each relaxing all $m$
edges in $\Th(m)$ time, plus one detector pass. Total:
\[
\Th\big((n-1)\cdot m + m\big) = \Th(nm).
\]
With $m$ up to $\Th(n^2)$ in a dense graph, that is $\Th(n^3)$ in the worst case.
Space is $\Th(n)$ for the $d$ and $\pi$ arrays. It is slow, but it is the price of
handling negative edges and detecting negative cycles.
\end{complexitybox}

\begin{trapbox}
A common implementation bug: people ``optimize'' by stopping early when a pass
makes no changes (a fine and correct optimization) but then \emph{forget the
detector pass}, or they run only $n-2$ passes off a fencepost error. If the longest
shortest path has exactly $n-1$ edges, $n-2$ passes leave one vertex wrong. Always
$n-1$ relaxation passes, then one detector pass. Count carefully.
\end{trapbox}

\section{Shortest paths in a DAG: order is free}

Now a special case where we can be genuinely clever about order. Suppose the graph
is a \emph{directed acyclic graph} (DAG) -- no cycles at all. Then there are no
negative cycles to worry about (no cycles, period), and something wonderful is
available to us: a \emph{topological order}.

Recall from Chapter~22 that a topological sort lists the vertices so that every
edge $(u,v)$ goes from an earlier vertex to a later one. In a DAG such an ordering
always exists and is computable in $\Th(n+m)$ time. Here is why it is exactly what
the convergence lemma craves.

\begin{intuitionbox}
Imagine the vertices lined up left to right in topological order, every edge
pointing rightward. Process them left to right. When you arrive at a vertex $u$,
\emph{every} vertex that could possibly feed into $u$ -- everything that could
lower $d[u]$ -- sits to its left and has already been processed. So $d[u]$ is
already final the moment you reach $u$. You relax $u$'s outgoing edges once, push
its now-correct estimate rightward, and never look back. One pass. No repetition.
No priority queue. Just a single left-to-right sweep
(Figure~\ref{ch24:dag-sweep}).
\end{intuitionbox}

\begin{figure}[h]
\centering
\begin{tikzpicture}[
  scale=0.92,
  every node/.style={font=\small},
  vert/.style={circle,draw=gray!70,thick,minimum size=8mm,inner sep=0pt,font=\small},
  edgelabel/.style={fill=white,inner sep=1pt,font=\tiny},
  lbl/.style={font=\scriptsize\scshape, text=gray!80!black}
]
  % Top Lane: Topological DAG with Edge Weights
  \draw[rounded corners=4pt, fill=gray!3, draw=gray!25] (-1.2, 0.4) rectangle (9.6, 3.8);
  \node[lbl, anchor=west] at (-1.0, 3.5) {Topological Order: All Directed Edges Flow Forward};

  \node[vert, fill=sourceblue!25, draw=sourceblue!80!black] (s) at (0, 1.8) {$s$};
  \node[vert, fill=white] (a) at (2.0, 2.7) {$a$};
  \node[vert, fill=white] (b) at (4.0, 1.1) {$b$};
  \node[vert, fill=white] (c) at (6.0, 2.7) {$c$};
  \node[vert, fill=white] (d) at (8.4, 1.8) {$d$};

  \draw[->, >=stealth, thick] (s) -- node[edgelabel, midway] {$2$} (a);
  \draw[->, >=stealth, thick] (s) to[bend right=12] node[edgelabel, midway] {$6$} (b);
  \draw[->, >=stealth, thick] (a) -- node[edgelabel, midway] {$3$} (b);
  \draw[->, >=stealth, thick] (a) to[bend left=16] node[edgelabel, midway] {$5$} (c);
  \draw[->, >=stealth, thick] (b) -- node[edgelabel, midway] {$1$} (c);
  \draw[->, >=stealth, thick] (b) to[bend right=16] node[edgelabel, midway] {$8$} (d);
  \draw[->, >=stealth, thick] (c) -- node[edgelabel, midway] {$4$} (d);

  % Bottom Lane: Sequential Sweep Progression Table
  \draw[rounded corners=4pt, fill=sourceblue!2, draw=sourceblue!25] (-1.2, -3.3) rectangle (9.6, -0.1);
  \node[lbl, anchor=west] at (-1.0, -0.4) {Single-Pass Sweep: Distance Estimates Freeze in Topological Order};

  \draw[->, structurepurple, very thick] (-0.8, -0.85) -- (9.0, -0.85)
    node[midway, below, font=\scriptsize\bfseries, text=structurepurple!90!black] {left-to-right processing sweep $\longrightarrow$};

  \node[font=\tiny\ttfamily, anchor=west] at (-0.9, -1.35) {At s: d[s]=0 final  $\to$ relax (s,a): d[a]=2, relax (s,b): d[b]=6};
  \node[font=\tiny\ttfamily, anchor=west] at (-0.9, -1.75) {At a: d[a]=2 final  $\to$ relax (a,b): d[b]=min(6, 2+3)=5, relax (a,c): d[c]=7};
  \node[font=\tiny\ttfamily, anchor=west] at (-0.9, -2.15) {At b: d[b]=5 final  $\to$ relax (b,c): d[c]=min(7, 5+1)=6, relax (b,d): d[d]=13};
  \node[font=\tiny\ttfamily, anchor=west] at (-0.9, -2.55) {At c: d[c]=6 final  $\to$ relax (c,d): d[d]=min(13, 6+4)=10};
  \node[font=\tiny\bfseries, text=targetgreen!80!black, anchor=west] at (-0.9, -2.95) {At d: d[d]=10 final $\to$ All shortest paths computed in $\Theta(V+E)$ total time.};
\end{tikzpicture}
\caption{A directed acyclic graph processed in topological order. Top: every edge points strictly rightward, ensuring that when the sweep reaches vertex $u$, all potential predecessors have already transmitted their final distances. Bottom: step-by-step relaxation timeline showing how a single linear sweep locks in optimal distances without any priority queue overhead.}
\label{ch24:dag-sweep}
\end{figure}

\begin{algorithm}[H]
\caption{\textsc{DAG-Shortest-Paths}($G, w, s$)}
\KwIn{Weighted DAG $G=(V,E)$, weights $w$ (any sign), source $s$.}
\KwOut{Correct $d$ and $\pi$.}
topologically sort the vertices of $G$\;
\textsc{Initialize-Single-Source}($G,s$)\;
\For{each vertex $u$ in topological order}{
  \For{each vertex $v$ in $\mathrm{Adj}[u]$}{
    \textsc{Relax}($u, v, w$)\;
  }
}
\end{algorithm}

\begin{formalbox}
\emph{Why it is correct.} We claim $d[u]=\delta(s,u)$ for every $u$, by induction
along the topological order. Consider a shortest path
$s = v_0 \to v_1 \to \dots \to v_k = u$. Because every edge goes forward in
topological order, the vertices $v_0, v_1, \dots, v_k$ appear in increasing
topological position. The algorithm processes them in that very order, and when it
processes $v_{i-1}$ it relaxes $(v_{i-1}, v_i)$. By induction $d[v_{i-1}]$ is
already $\delta(s,v_{i-1})$ at that moment, so by the convergence lemma
(Lemma~\ref{ch24:lemma-convergence}) $d[v_i]$ becomes $\delta(s,v_i)$. Marching
$i$ from $1$ to $k$ gives $d[u]=\delta(s,u)$. \qed
\end{formalbox}

\begin{complexitybox}
\textbf{DAG shortest paths.} Topological sort is $\Th(n+m)$. The main loop touches
each vertex once and each edge exactly once (every edge is relaxed a single time),
which is again $\Th(n+m)$. Total: $\Th(n+m)$ -- \emph{linear}. This is the fastest
single-source algorithm in the chapter, and it works with negative edges, because
a DAG can never contain a negative cycle.
\end{complexitybox}

\begin{aestheticbox}
DAG shortest paths is the engine behind \emph{critical-path analysis} in project
scheduling (PERT charts). Tasks are vertices, dependencies are edges, durations
are weights, and the \emph{longest} path -- found by negating weights and running
this same algorithm -- is the minimum time to finish the whole project. The same
linear sweep also underlies many dynamic programs: a DP over a DAG of subproblems
is literally this algorithm in disguise, which is why Chapter~15's tables and this
chapter's relaxations feel so similar.
\end{aestheticbox}

\section{Dijkstra's algorithm: greedy on nonnegative weights}

Now the famous one. Edsger Dijkstra's algorithm assumes \emph{all edge weights are
nonnegative} and, in return, is fast and elegant. It does not need a DAG, and it
does not relax every edge $n-1$ times. Instead it discovers the correct order on
the fly using a greedy rule.

\begin{intuitionbox}
Think of it as fire spreading from the source, but a careful, accountant's fire.
Each vertex has a current best-known time $d[v]$. Maintain a set $S$ of vertices
whose true distance is already \emph{finalized}. Repeatedly do this: among all
vertices \emph{not} yet finalized, pick the one with the smallest current estimate
$d[u]$ -- call it the ``closest frontier vertex'' -- declare its estimate final,
add it to $S$, and relax all its outgoing edges to update its neighbors. Repeat
until every vertex is finalized. The closest unfinished vertex is always safe to
finalize, because (with nonnegative weights) no cheaper route to it could possibly
exist through a vertex that is even farther away.
\end{intuitionbox}

That last sentence is the whole algorithm, and also exactly where nonnegativity is
load-bearing. Hold it; we will prove it.

\begin{algorithm}[H]
\caption{\textsc{Dijkstra}($G, w, s$)}
\KwIn{Weighted digraph $G=(V,E)$ with $w(u,v)\ge 0$ for all edges, source $s$.}
\KwOut{Correct $d$ and $\pi$.}
\textsc{Initialize-Single-Source}($G,s$)\;
$S \leftarrow \emptyset$\;
$Q \leftarrow V$ \tcp*{min-priority queue keyed on $d$}
\While{$Q \ne \emptyset$}{
  $u \leftarrow \textsc{Extract-Min}(Q)$ \tcp*{closest unfinalized vertex}
  $S \leftarrow S \cup \{u\}$\;
  \For{each vertex $v \in \mathrm{Adj}[u]$}{
    \textsc{Relax}($u, v, w$) \tcp*{may call \textsc{Decrease-Key} on $v$}
  }
}
\end{algorithm}

Each vertex is extracted from $Q$ exactly once, and at the moment of extraction we
claim its estimate is already the true distance. The greedy choice -- always pull
the smallest -- is what makes the order come out right.

\subsection{A complete numeric trace}

Let us run Dijkstra on a small \emph{nonnegative} graph. Vertices $s,t,x,y,z$,
edges:
\[
\begin{aligned}
&(s,t)=10,\ (s,y)=5,\ (t,x)=1,\ (t,y)=2,\\
&(y,t)=3,\ (y,x)=9,\ (y,z)=2,\ (x,z)=4,\ (z,x)=6,\ (z,s)=7.
\end{aligned}
\]

\begin{figure}[h]
\centering
\begin{tikzpicture}[
  scale=0.92,
  every node/.style={font=\small},
  vert/.style={circle,draw=gray!70,thick,minimum size=7mm,inner sep=0pt,font=\small},
  edgelabel/.style={fill=white,inner sep=1pt,font=\tiny},
  lbl/.style={font=\scriptsize\scshape, text=gray!80!black},
  qbox/.style={draw=gray!50, rounded corners=2pt, fill=cyan!8, minimum width=1.1cm, minimum height=0.6cm, font=\scriptsize}
]
  % Panel A: Priority Queue State & Extraction
  \draw[rounded corners=4pt, fill=gray!3, draw=gray!25] (-0.8, -1.2) rectangle (5.2, 4.2);
  \node[lbl, anchor=west] at (-0.6, 3.9) {Priority Queue Extraction};

  \node[font=\scriptsize\bfseries, text=cyan!80!black] at (2.2, 3.1) {Min-Heap $Q$ before extraction:};
  \node[qbox, fill=quotientorange!25, draw=quotientorange!80!black] (qy) at (0.2, 2.3) {$y: 5$};
  \node[qbox] (qt) at (1.5, 2.3) {$t: 10$};
  \node[qbox] (qz) at (2.8, 2.3) {$z: \infty$};
  \node[qbox] (qx) at (4.1, 2.3) {$x: \infty$};

  \draw[->, >=stealth, very thick, quotientorange!90!black] (qy.south) -- ++(0, -0.6)
    node[midway, right, font=\tiny\bfseries] {Extract-Min};

  \node[vert, fill=quotientorange!30, draw=quotientorange!90!black] (extr_y) at (0.2, 0.7) {$y$};
  \node[font=\tiny, text=quotientorange!90!black, below=2pt of extr_y] {Key $d[y]=5$ finalized};

  \draw[rounded corners=3pt, fill=sourceblue!10, draw=sourceblue!40] (1.8, 0.2) rectangle (4.8, 1.2);
  \node[font=\tiny\bfseries, text=sourceblue!90!black] at (3.3, 0.85) {Settled Set $S \leftarrow S \cup \{y\}$};
  \node[font=\tiny, text=gray!70!black] at (3.3, 0.45) {$\delta(s,y)=5$ proven optimal};

  % Panel B: Relaxing Outgoing Edges & Decrease-Key
  \begin{scope}[xshift=6.4cm]
    \draw[rounded corners=4pt, fill=gray!3, draw=gray!25] (-0.8, -1.2) rectangle (6.2, 4.2);
    \node[lbl, anchor=west] at (-0.6, 3.9) {Relaxing Outgoing Edges of $y$};

    \node[vert, fill=quotientorange!30, draw=quotientorange!90!black] (py) at (0.4, 1.6) {$y$};
    \node[vert, fill=white] (pt) at (3.2, 2.7) {$t$};
    \node[vert, fill=white] (pz) at (5.0, 1.6) {$z$};
    \node[vert, fill=white] (px) at (3.2, 0.2) {$x$};

    % Relaxations with labels
    \draw[->, >=stealth, very thick, trapred] (py) -- node[edgelabel, midway] {$w=3$} (pt);
    \draw[->, >=stealth, very thick, trapred] (py) -- node[edgelabel, midway] {$w=2$} (pz);
    \draw[->, >=stealth, very thick, trapred] (py) -- node[edgelabel, midway] {$w=9$} (px);

    % Post-relaxation keys placed compactly by each vertex
    \node[font=\tiny\bfseries, text=trapred!90!black, above=2pt of pt] {$d[t]: 10 \to 8$};
    \node[font=\tiny\bfseries, text=trapred!90!black, right=3pt of pz] {$d[z]: \infty \to 7$};
    \node[font=\tiny\bfseries, text=trapred!90!black, below=2pt of px] {$d[x]: \infty \to 14$};

    \node[font=\scriptsize\bfseries, text=targetgreen!80!black] at (2.7, -0.85) {Updated $Q$: $[z{:}7,\ t{:}8,\ x{:}14]$};
  \end{scope}
\end{tikzpicture}
\caption{Dijkstra priority queue extraction versus outgoing edge relaxation. Left: the min-priority queue $Q$ yields the vertex with smallest tentative distance estimate, $y$ with $d[y]=5$, which moves into the settled set $S$. Right: the outgoing edges of $y$ fire relaxation updates to its neighbors, triggering \textsc{Decrease-Key} operations in $Q$ for $t$ (lowered to $8$), $z$ (lowered to $7$), and $x$ (lowered to $14$).}
\label{ch24:dij-graph}
\end{figure}

\begin{labbox}
Initialize: $d[s]=0$, rest $\infty$; $S=\emptyset$. At each step I pick the
non-finalized vertex of smallest $d$, finalize it, and relax its out-edges.

\medskip
\textbf{Step 1.} Smallest $d$ is $s$ ($d=0$). Finalize $s$. Relax: $(s,t)$ gives
$d[t]=10$; $(s,y)$ gives $d[y]=5$.
\[
S=\{s\},\quad d=(t{:}10,\ y{:}5,\ x{:}\infty,\ z{:}\infty).
\]

\medskip
\textbf{Step 2.} Smallest non-final is $y$ ($d=5$). Finalize $y$. Relax:
$(y,t)$: $5+3=8 < 10$, so $d[t]=8$. $(y,x)$: $5+9=14$, so $d[x]=14$. $(y,z)$:
$5+2=7$, so $d[z]=7$.
\[
S=\{s,y\},\quad d=(t{:}8,\ x{:}14,\ z{:}7).
\]

\medskip
\textbf{Step 3.} Smallest non-final is $z$ ($d=7$). Finalize $z$. Relax:
$(z,x)$: $7+6=13 < 14$, so $d[x]=13$. $(z,s)$: $7+7=14$, but $s$ is finalized at
$0$ -- no change.
\[
S=\{s,y,z\},\quad d=(t{:}8,\ x{:}13).
\]

\medskip
\textbf{Step 4.} Smallest non-final is $t$ ($d=8$). Finalize $t$. Relax:
$(t,x)$: $8+1=9 < 13$, so $d[x]=9$. $(t,y)$: $y$ already finalized -- no change.
\[
S=\{s,y,z,t\},\quad d=(x{:}9).
\]

\medskip
\textbf{Step 5.} Only $x$ remains ($d=9$). Finalize it. Done.
\[
\delta(s,s)=0,\ \delta(s,y)=5,\ \delta(s,z)=7,\ \delta(s,t)=8,\ \delta(s,x)=9.
\]
\end{labbox}

Notice the finalization order: $s, y, z, t, x$ -- exactly increasing $d$-value,
$0,5,7,8,9$. That monotonicity holds in general. Dijkstra finalizes vertices in nondecreasing
order of true distance. That monotonic order is the soul of the correctness proof.

\subsection{Why Dijkstra is correct}

\begin{theorem}[Dijkstra's correctness]\label{ch24:thm-dijkstra}
Run on a graph with all nonnegative edge weights, \textsc{Dijkstra} terminates
with $d[u]=\delta(s,u)$ for every vertex $u$.
\end{theorem}

The key is the invariant that the vertex we extract is already correct. Let us
prove it carefully but gently. The picture to hold in mind is the cut between the
finalized set $S$ and the rest (Figure~\ref{ch24:dij-cut}): a shortest path to the
soon-to-be-extracted $u$ must cross that cut somewhere, and we zero in on the very
first crossing edge.

\begin{figure}[h]
\centering
\begin{tikzpicture}[>=Stealth,
  every node/.style={circle,draw,minimum size=7mm,inner sep=1pt,font=\small}]
  \begin{scope}[on background layer]
    \fill[sourceblue!12,rounded corners] (-0.7,-1.5) rectangle (3.2,1.5);
  \end{scope}
  \node[draw=none,fill=none,font=\scriptsize\bfseries,sourceblue] at (1.2,1.85) {finalized set $S$};
  \node[fill=sourceblue!25] (s) at (0,0.7) {$s$};
  \node (mid) at (1.6,-0.5) {};
  \node (x) at (2.6,0.4) {$x$};
  \node (y) at (5,0.4) {$y$};
  \node (u) at (7,-0.4) {$u$};
  \draw[->,dashed] (s) to[bend right=10] (mid);
  \draw[->,dashed] (mid) to[bend right=10] (x);
  \draw[->,trapred,very thick] (x) -- node[above,draw=none,fill=none,font=\scriptsize]{first edge leaving $S$} (y);
  \draw[->,dashed] (y) to[bend left=12] (u);
  \node[draw=none,fill=none,font=\scriptsize] at (5,-1.0) {$d[y]=\delta(s,y)$};
  \node[draw=none,fill=none,font=\scriptsize] at (7,-1.2) {extracting $u$};
\end{tikzpicture}
\caption{The cut argument in Dijkstra's proof. A shortest path from $s$ (inside
$S$) to $u$ (outside $S$) crosses the boundary; $(x,y)$ is its first crossing edge.
We show $d[y]$ is already correct and $d[y]\le d[u]$, forcing $d[u]=\delta(s,u)$.}
\label{ch24:dij-cut}
\end{figure}

\begin{formalbox}
\emph{Proof.} We show: \emph{at the moment a vertex $u$ is extracted (added to
$S$), $d[u]=\delta(s,u)$.} Since $d$ never changes after extraction
(Lemma~\ref{ch24:lemma-upperbound}, $d[u]$ already at its floor), this gives the
theorem.

Suppose not. Let $u$ be the \emph{first} vertex extracted with
$d[u]\ne\delta(s,u)$; by Lemma~\ref{ch24:lemma-upperbound} this means
$d[u]>\delta(s,u)$, so in particular $u\ne s$ (since $s$ is extracted first with
$d[s]=0=\delta(s,s)$) and $u$ is reachable (else $\delta=\infty=d$).

Consider a shortest path $p$ from $s$ to $u$. It starts inside $S$ (at $s$) and
ends outside $S$ (at $u$, the vertex we are about to add). Walk along $p$ and let
$(x,y)$ be the \emph{first} edge that leaves $S$: $x\in S$, $y\notin S$, with $y$
the first off-$S$ vertex on $p$ (possibly $y=u$).

\smallskip
Because $x$ was extracted before $u$ and $u$ is the \emph{first} bad extraction,
$x$ was correct when extracted: $d[x]=\delta(s,x)$. When $x$ was extracted we
relaxed the edge $(x,y)$, so by the convergence lemma
(Lemma~\ref{ch24:lemma-convergence}) $d[y]=\delta(s,y)$ now.

\smallskip
Now the punchline, and the place nonnegativity is essential. Since $y$ lies on a
shortest path to $u$ and all weights from $y$ onward to $u$ are
\emph{nonnegative}, the subpath from $y$ to $u$ has weight $\ge 0$, so
$\delta(s,y) \le \delta(s,u)$. Therefore
\[
d[y] = \delta(s,y) \le \delta(s,u) \le d[u],
\]
the last step because $d[u]\ge\delta(s,u)$ always. But $u$ was chosen by
\textsc{Extract-Min} as having the \emph{smallest} $d$ among non-$S$ vertices, and
$y$ is also outside $S$, so $d[u]\le d[y]$. Squeezing, $d[y]=d[u]=\delta(s,u)$ --
which says $d[u]=\delta(s,u)$, contradicting our assumption that $u$ was bad. No
bad extraction exists. \qed
\end{formalbox}

\begin{trapbox}
Stare at the step ``$\delta(s,y)\le\delta(s,u)$ because the subpath from $y$ to
$u$ has nonnegative weight.'' If a \emph{negative} edge lay on that subpath, the
inequality could flip: a vertex $y$ farther out by current estimate might
secretly reach $u$ via a negative shortcut for less. Then finalizing the smallest
$d$ is wrong, and Dijkstra returns incorrect answers \emph{silently}. This is the
precise, single reason Dijkstra requires nonnegative weights. Our earlier
Bellman--Ford trace, where the best path to $z$ snaked through negatives, is a
concrete graph on which Dijkstra would have failed.
\end{trapbox}

\subsection{The role of the priority queue}

Dijkstra's skeleton -- ``repeatedly grab the closest unfinalized vertex, finalize
it, relax its edges'' -- demands a data structure that supports three operations:
build the queue, \textsc{Extract-Min} (remove the smallest-$d$ vertex), and
\textsc{Decrease-Key} (lower a vertex's $d$ when a relaxation improves it). The
\emph{entire} running time of Dijkstra is determined by how fast that priority
queue is. The algorithm performs exactly $n$ \textsc{Extract-Min} operations (one
per vertex) and at most $m$ \textsc{Decrease-Key} operations (at most one per
edge). So the cost is
\[
n \cdot T_{\text{extract}} \;+\; m \cdot T_{\text{decrease}} \;+\; T_{\text{build}}.
\]
Different queues give different bounds:

\begin{itemize}
\item \textbf{Unsorted array / linear scan.} \textsc{Extract-Min} scans all
remaining vertices: $\Th(n)$ each. \textsc{Decrease-Key} is $\Th(1)$ (just write
the value). Total $\Th(n^2 + m) = \Th(n^2)$. Best for \emph{dense} graphs where
$m \approx n^2$.
\item \textbf{Binary min-heap} (Chapter~6). \textsc{Extract-Min} and
\textsc{Decrease-Key} are each $\Th(\log n)$. Total
$\Th((n+m)\log n) = \Th(m \log n)$ for a connected graph. Best for \emph{sparse}
graphs where $m = \Oh(n^2/\log n)$. This is the version most libraries ship.
\item \textbf{Fibonacci heap} (Chapter~20). \textsc{Extract-Min} is
$\Oh(\log n)$ amortized, but \textsc{Decrease-Key} is $\Oh(1)$ amortized. Total
$\Th(n\log n + m)$, the theoretical best. The constants are large, so it wins only
on very large, edge-heavy graphs.
\end{itemize}

\begin{complexitybox}
\textbf{Dijkstra, by priority queue.}
\[
\begin{array}{lll}
\text{array} & \Th(n^2) & \text{dense graphs}\\
\text{binary heap} & \Th(m\log n) & \text{sparse graphs, the practical default}\\
\text{Fibonacci heap} & \Th(n\log n + m) & \text{asymptotically optimal}
\end{array}
\]
Across all three the count of operations is identical -- $n$ extractions and $\le
m$ decrease-keys -- only the per-operation cost differs. This is a textbook example
of an algorithm whose performance is dictated by its data structure, not its
control flow.
\end{complexitybox}

\begin{historybox}
Edsger W.\ Dijkstra devised this algorithm in 1956 (published 1959), reportedly in
about twenty minutes at a cafe, while thinking about how to demonstrate the power
of a new computer to a general audience. Richard Bellman and Lester Ford developed
the relaxation-everything approach in the late 1950s. The unifying view -- that all
of these are relaxation in different orders -- is the modern textbook framing.
\emph{In CLRS this material is Chapter~22 (Single-Source Shortest Paths in the
4th edition; Chapter~24 in the 3rd): \S22.1 properties and relaxation, \S22.1
Bellman--Ford, \S22.2 DAG shortest paths, \S22.3 Dijkstra.} The all-pairs problem
and Johnson's algorithm (which uses Bellman--Ford to reweight a graph so Dijkstra
can run safely with no negative edges) come in the next chapter.
\end{historybox}

\subsection{The predecessor graph: reconstructing the actual path}

We have been computing \emph{weights} $d[v]$, but the user usually wants the
\emph{route}. That is what $\pi$ is for. After any of these algorithms finishes,
the $\pi$ pointers form a \emph{shortest-paths tree} rooted at $s$: follow $\pi$
backward from any vertex $v$ and you walk a shortest path from $v$ to $s$ in
reverse. Reversing that list gives the path from $s$ to $v$.

\begin{algorithm}[H]
\caption{\textsc{Print-Path}($G, s, v$)}
\KwIn{Source $s$, target $v$, predecessor array $\pi$.}
\KwOut{The vertices of a shortest $s\!\to\!v$ path, in order.}
\If{$v = s$}{
  print $s$\;
}
\ElseIf{$\pi[v] = \textsc{nil}$}{
  print ``no path from $s$ to $v$''\;
}
\Else{
  \textsc{Print-Path}($G, s, \pi[v]$)\;
  print $v$\;
}
\end{algorithm}

\begin{focusbox}
The shape of a shortest-paths tree: $s$ is the root; each vertex $v$ has exactly
one parent $\pi[v]$; and the tree path from $s$ to $v$ is \emph{a} shortest path.
There may be other shortest paths of equal weight, but the tree commits to one. The
relaxation steps build this tree edge by edge: every time \textsc{Relax} updates
$\pi[v]\leftarrow u$, it is re-parenting $v$ under a better-positioned $u$.
\end{focusbox}

\section{Putting the three together}

Three algorithms, one operation. Here is the map to keep in your head.

\begin{center}
\begin{tabular}{@{}lllll@{}}
\toprule
Algorithm & Graph & Negative edges? & Order of relaxation & Time \\
\midrule
Bellman--Ford & any & yes (detects neg.\ cycles) & all edges, $n-1$ times & $\Th(nm)$ \\
DAG & acyclic only & yes & topological order, once & $\Th(n+m)$ \\
Dijkstra & any & \emph{no} & greedy, closest-first & $\Th(m\log n)$ \\
\bottomrule
\end{tabular}
\end{center}

\begin{intuitionbox}
A decision tree for choosing. \emph{Is the graph a DAG?} If yes, use the linear
DAG algorithm -- nothing beats $\Th(n+m)$. \emph{Are there negative edges (and the
graph is not a DAG)?} You must use Bellman--Ford, and you get negative-cycle
detection for free. \emph{Otherwise} -- general graph, nonnegative weights, the
overwhelmingly common case -- use Dijkstra with a binary heap. Faster, cleaner, and
the basis of real routing engines (which then add tricks like A* and contraction
hierarchies on top).
\end{intuitionbox}

\begin{aestheticbox}
These three power a startling range of real systems. Internet routing protocols
(OSPF) run Dijkstra to compute forwarding tables. GPS navigation runs Dijkstra
variants over road graphs. Currency-arbitrage detection runs Bellman--Ford: take
logs of exchange rates, negate, and a reachable negative cycle is a sequence of
trades that prints money. Build systems and spreadsheet recalculation run the DAG
algorithm over dependency graphs. The same five-line \textsc{Relax} sits at the
heart of all of them.
\end{aestheticbox}

\begin{recoverybox}
If this chapter felt like a lot, here is the one-paragraph version to carry away.
A shortest-path estimate $d[v]$ starts at $\infty$ and only ever decreases.
\emph{Relaxing} an edge $(u,v)$ checks whether routing through $u$ beats $v$'s
current estimate, and if so improves it. Relax the edges of a shortest path in
order and its endpoint's estimate becomes exactly correct. The three algorithms
are three disciplines for guaranteeing good relaxation order: brute repetition
(Bellman--Ford), topological order (DAG), and greedy closest-first
(Dijkstra). Get those, and you have the chapter.
\end{recoverybox}

\section{Exercises}

\begin{exercise}\label{ch24:ex-relax-order}
Consider the graph with vertices $s,a,b,c$ and edges $(s,a)=4$, $(s,b)=1$,
$(b,a)=2$, $(a,c)=1$, $(b,c)=5$. Initialize and then relax the edges in the order
$(s,a)$, $(a,c)$, $(b,c)$, $(s,b)$, $(b,a)$, $(a,c)$. Show $d$ after each relaxation. Did $d[c]$
reach $\delta(s,c)$? Explain how the order helped or hurt.
\end{exercise}

\begin{exercise}\label{ch24:ex-bf-trace}
Run Bellman--Ford by hand on the graph $s,a,b,c$ with edges $(s,a)=3$, $(s,b)=5$,
$(a,b)=-2$, $(b,c)=2$, $(a,c)=6$, $(c,a)=1$. Use any fixed edge order, show the
$d$-array after each of the $n-1$ passes, and state whether the detector pass
returns \textsc{true} or \textsc{false}. (Check first: the cycle
$a\to b\to c\to a$ has weight $-2+2+1=1>0$, so there is no negative cycle and the
detector should return \textsc{true}.)
\end{exercise}

\begin{exercise}\label{ch24:ex-neg-cycle}
Modify edge $(c,a)$ in the previous exercise to weight $-6$. Now the cycle
$a\to b\to c\to a$ becomes a negative-weight cycle reachable from $s$. Identify it,
compute its total weight, and confirm that the Bellman--Ford detector would return
\textsc{false} (no $n-1$ passes will settle the $d$-values, and the detector pass
finds an edge it can still relax).
\end{exercise}

\begin{exercise}\label{ch24:ex-dag}
Give a topological order for the DAG with edges
$(s,a),(s,c),(a,b),(c,a),(c,d),(b,d)$ and weights
$2,4,1,3,1,2$ respectively. Run \textsc{DAG-Shortest-Paths} from $s$ and report
$d$ and $\pi$ for every vertex. How many times is each edge relaxed?
\end{exercise}

\begin{exercise}\label{ch24:ex-dijkstra-trace}
Run Dijkstra from $s$ on the graph with edges $(s,a)=2$, $(s,b)=6$, $(a,b)=3$,
$(a,c)=5$, $(b,c)=1$, $(c,d)=4$, $(b,d)=8$ (all nonnegative). List the order in
which vertices are finalized and the final $d$-values. Verify that the
finalization order is nondecreasing in $d$.
\end{exercise}

\begin{exercise}\label{ch24:ex-dijkstra-fail}
Construct a graph with \emph{one} negative edge and no negative cycle on which
Dijkstra returns a wrong distance for some vertex. Trace Dijkstra to show the
wrong answer, then run Bellman--Ford to show the right one. (Three vertices
suffice.)
\end{exercise}

\begin{exercise}\label{ch24:ex-pq-counts}
Argue carefully that Dijkstra performs exactly $n$ \textsc{Extract-Min} operations
and at most $m$ \textsc{Decrease-Key} operations, regardless of which priority
queue is used. Now specialize to a $d$-ary heap (a heap where each node has $d$
children), where \textsc{Extract-Min} costs $\Th\!\big(d\,\frac{\log n}{\log d}\big)$
and \textsc{Decrease-Key} costs $\Th\!\big(\frac{\log n}{\log d}\big)$. Write the
total running time as a function of $n$, $m$, and $d$, and show that choosing
$d \approx \max(2,\, m/n)$ minimizes it. What total time does that choice give for a
dense graph with $m = \Th(n^2)$? (Binary heap, $d=2$, and the dense-graph optimum
should both fall out of your formula.)
\end{exercise}

\begin{exercise}\label{ch24:ex-longest-dag}
You are given a DAG of tasks with positive durations on the edges and want the
\emph{longest} path (the critical path / minimum project completion time). Explain
precisely how to adapt \textsc{DAG-Shortest-Paths} to compute longest paths, and
why the same trick would \emph{not} work for longest paths in a general graph with
cycles.
\end{exercise}

\begin{exercise}\label{ch24:ex-invariant}
Prove the upper-bound invariant (Lemma~\ref{ch24:lemma-upperbound}) in full,
filling in every inductive step, including the base case and the case where a
relaxation does \emph{not} change $d[v]$. Where exactly does the triangle
inequality get used?
\end{exercise}

\begin{exercise}\label{ch24:ex-why-n-minus-1}
Explain, in your own words and with a worst-case example graph, why Bellman--Ford
needs $n-1$ passes and not fewer. Construct a graph on $n=5$ vertices where some
vertex's $d$-value only becomes correct on the very last (fourth) pass. (Hint: a
single long path $s\to v_1\to v_2\to v_3\to v_4$, and pick an edge order that
relaxes the path edges in reverse each pass.)
\end{exercise}
